\chapter{Conclusioni e sviluppi futuri}

Il progetto ZTALeaks ha esplorato l'applicazione dei principi del modello Zero Trust (NIST SP 800-207) attraverso la progettazione e l'implementazione di un'architettura di riferimento containerizzata. L'adozione del caso d'uso di una centrale nucleare, caratterizzato da severi requisiti di confidenzialità e ruoli operativi ben definiti, ha permesso di validare le scelte architetturali e i meccanismi di controllo degli accessi in uno scenario ad alto rischio, distribuendo le responsabilità di sicurezza tra i vari componenti del sistema.

\section{Risultati ottenuti}

Dal punto di vista architetturale, l'implementazione ha mantenuto una separazione rigorosa tra Policy Enforcement Point (Envoy Proxy) e Policy Decision Point (Security Orchestrator, OPA e modello AI). Questa suddivisione, unita a una topologia di rete articolata in segmenti Docker isolati (\texttt{Front-Net}, \texttt{Auth-Net}, \texttt{Back-Net} e \texttt{Snort-Net}), ha consentito di applicare il principio della \textit{complete mediation} a livello di rete, vincolando il traffico ai soli percorsi autorizzati.

Il controllo degli accessi è stato strutturato su tre livelli complementari: un livello Role-Based Access Control (RBAC), per mappare i ruoli aziendali estratti dai token JWT in etichette di sicurezza; un livello Attribute-Based Access Control (ABAC), per la valutazione contestuale delle richieste (percorso, metodo HTTP, dispositivo e score di rischio dinamico); e un'implementazione del modello Bell-LaPadula tramite policy Rego in OPA, per imporre formalmente la \textit{Simple Security Property} e la \textit{Star Property}, includendo il meccanismo del \textit{Trusted Subject} per la declassificazione controllata.

L'introduzione di un algoritmo di fiducia ibrido ha permesso di estendere i controlli statici con una valutazione continua del rischio fornita dal modello Graphagate. Le decisioni di accesso vengono prese ponderando il beneficio dell'operazione con lo score di anomalia contestuale, richiedendo che il beneficio netto superi una soglia predefinita per ciascun endpoint. Tale approccio si è dimostrato efficace nel mitigare potenziali attacchi condotti con credenziali valide ma associati a pattern comportamentali anomali.

Per quanto riguarda la rilevazione delle minacce, il modello Temporal Graph Network (Graphagate) ha ottenuto una AUC aggregata di $0.965 \pm 0.007$ e una recall aggregata pari a $0.811$ sul set di test sintetico, e una AUC specifica sul movimento laterale pari a $0.922$, dimostrando prestazioni di rilievo rispetto alle baseline non relazionali testate. L'evoluzione dello schema dati alla versione v4, con l'introduzione del nodo \textit{Configuration}, ha migliorato la capacità di rilevazione del movimento laterale e permesso l'identificazione di simulazioni di furto di credenziali. È stato inoltre validato il meccanismo anti-poisoning, che previene l'inquinamento del modello condizionando l'aggiornamento dello stato comportamentale all'esito positivo delle valutazioni in OPA.

L'infrastruttura di osservabilità, basata sulla propagazione dell'\texttt{X-Request-ID} e sulla centralizzazione dei log in Splunk, ha garantito una tracciabilità completa delle richieste attraverso  firewall, proxy, orchestratore e servizi di backend, integrata da un audit a livello di base dati che registra ogni operazione MongoDB associandola al microservizio e all'utente finale responsabili. I test di carico hanno evidenziato tempi di risposta medi tra i 16 e i 32 ms per carichi fino a 20 utenti concorrenti, con una latenza del modello TGN adeguata per il percorso decisionale sincrono ($P50 \approx 9.6$ ms). Le simulazioni di attacco hanno infine confermato la capacità della pipeline di rilevare e bloccare varie tipologie di minacce, tra cui SQL injection, XSS, connessioni TLS non sicure, attacchi Slowloris e movimenti laterali.

\section{Limiti e sviluppi futuri}

L'architettura proposta, pur dimostrando l'applicabilità dei principi Zero Trust nel contesto in esame, presenta alcuni limiti architetturali e operativi che delineano diverse aree di miglioramento per lavori futuri.

Un primo limite riguarda  l'assenza di un meccanismo di revoca (come CRL o OCSP) che rappresenta una vulnerabilità in scenari ad alto rischio. Come sviluppo futuro, si può prevedere l'adozione di certificati a vita breve emessi da una Certificate Authority interna automatizzata, un approccio che ridurrebbe l'esposizione ai rischi di compromissione senza introdurre il carico infrastrutturale dei sistemi di revoca tradizionali.

Per quanto concerne il controllo degli accessi, l'attuale implementazione del modello Bell-LaPadula non copre la \textit{Discretionary Security Property} (DSP). La sua integrazione permetterebbe un controllo più granulare, sovrapponendo una matrice degli accessi discrezionale ai vincoli mandatori.

L'impiego del modello di machine learning presenta un vincolo di scalabilità orizzontale: lo stato comportamentale risiede nella RAM del singolo worker, rendendo complessa la replicazione senza degradare la coerenza del segnale temporale. Un possibile sviluppo futuro prevede l'introduzione di tecniche di sharding coerente per-entità; oppure si potrebbe utilizzare un database a grafo ad elevate prestazioni, così da permettere di avere più istanze di modello che condividono però la stessa memoria. Inoltre, per il consolidamento del modello, risulterebbe fondamentale validarne le prestazioni su dataset di traffico di rete reale, al fine di verificare le metriche di classificazione in ambienti operativi non controllati.

A livello infrastrutturale, la transizione dall'attuale orchestrazione basata su Docker Compose verso un cluster Kubernetes (per il quale sono state già approntate configurazioni preliminari) conferirebbe maggiore resilienza, auto-healing e scalabilità al sistema. In tale contesto, l'introduzione di un service mesh come Istio permetterebbe di estendere in modo nativo le garanzie del \textit{mutual TLS} alle comunicazioni interne tra i microservizi.

In conclusione, il progetto ZTALeaks offre un prototipo funzionale che integra modelli formali di controllo degli accessi, analisi comportamentale basata su Intelligenza Artificiale e tecniche di sicurezza infrastrutturale. Questo elaborato dimostra la fattibilità di un'architettura Zero Trust coerente e tracciabile per la protezione di un dominio critico, fornendo allo stesso tempo una base solida e verificabile per successive evoluzioni e applicazioni nel settore della sicurezza informatica.